% ============================================================
% 40_colors_structure.tex — Farbpaletten & Kapitel-Orchestrierung
% ============================================================

% ------------------------------------------------------------
% Index-Palette (18 Slots, kuratiert fuer 3-Fenster-Kontrast)
% ------------------------------------------------------------
% Reihenfolge: ETH-Signaturfarben zuerst, danach global distanzierte
% Zusatzfarben. Benachbarte Kapitel (Abstand 1 UND 2) bleiben deutlich
% unterscheidbar; globale Fast-Dopplungen werden vermieden, soweit es mit
% 18 kontraststarken Slots und fixierten ETH-Seeds sinnvoll moeglich ist.
% Alle Vollfarben sind dunkel genug fuer farbige Querverweise auf Weiss;
% die *Text-Tokens halten Balkentext austauschbar.
% Slot 0 ist reserviert fuer den Frontchapter (\StartFrontChapter) -- ETH-Blau.
% Niemals ChapterColor* in Kapiteldateien hart-coden; über \chaptercolor /
% \chaptercolorlight / \chaptercolortext referenzieren.
\definecolor{ChapterColor0}{HTML}{215CAF}       % ETH Blue
\definecolor{ChapterColor0Light}{HTML}{CEDBED}
\definecolor{ChapterColor0Text}{HTML}{FFFFFF}
\definecolor{ChapterColor1}{HTML}{6E3B6E}       % ETH Purple
\definecolor{ChapterColor1Light}{HTML}{DFD4DF}
\definecolor{ChapterColor1Text}{HTML}{FFFFFF}
\definecolor{ChapterColor2}{HTML}{5A7F2B}       % ETH Green
\definecolor{ChapterColor2Light}{HTML}{DBE3D0}
\definecolor{ChapterColor2Text}{HTML}{FFFFFF}
\definecolor{ChapterColor3}{HTML}{8C6239}       % ETH Bronze
\definecolor{ChapterColor3Light}{HTML}{E6DCD3}
\definecolor{ChapterColor3Text}{HTML}{FFFFFF}
\definecolor{ChapterColor4}{HTML}{007894}       % ETH Petrol
\definecolor{ChapterColor4Light}{HTML}{C7E1E7}
\definecolor{ChapterColor4Text}{HTML}{FFFFFF}
\definecolor{ChapterColor5}{HTML}{B7352D}       % ETH Red
\definecolor{ChapterColor5Light}{HTML}{EFD3D1}
\definecolor{ChapterColor5Text}{HTML}{FFFFFF}
\definecolor{ChapterColor6}{HTML}{23523B}       % Dark Spruce
\definecolor{ChapterColor6Light}{HTML}{CFD9D4}
\definecolor{ChapterColor6Text}{HTML}{FFFFFF}
\definecolor{ChapterColor7}{HTML}{8F2459}       % Raspberry
\definecolor{ChapterColor7Light}{HTML}{E6CFDA}
\definecolor{ChapterColor7Text}{HTML}{FFFFFF}
\definecolor{ChapterColor8}{HTML}{17205E}       % Midnight Blue
\definecolor{ChapterColor8Light}{HTML}{CCCEDC}
\definecolor{ChapterColor8Text}{HTML}{FFFFFF}
\definecolor{ChapterColor9}{HTML}{2D8659}       % Emerald
\definecolor{ChapterColor9Light}{HTML}{D1E4DA}
\definecolor{ChapterColor9Text}{HTML}{FFFFFF}
\definecolor{ChapterColor10}{HTML}{523523}      % Walnut
\definecolor{ChapterColor10Light}{HTML}{D9D3CF}
\definecolor{ChapterColor10Text}{HTML}{FFFFFF}
\definecolor{ChapterColor11}{HTML}{8F2481}      % Orchid
\definecolor{ChapterColor11Light}{HTML}{E6CFE3}
\definecolor{ChapterColor11Text}{HTML}{FFFFFF}
\definecolor{ChapterColor12}{HTML}{1D581D}      % Deep Green
\definecolor{ChapterColor12Light}{HTML}{CDDACD}
\definecolor{ChapterColor12Text}{HTML}{FFFFFF}
\definecolor{ChapterColor13}{HTML}{5E1720}      % Oxblood
\definecolor{ChapterColor13Light}{HTML}{DCCCCE}
\definecolor{ChapterColor13Text}{HTML}{FFFFFF}
\definecolor{ChapterColor14}{HTML}{4C248F}      % Violet
\definecolor{ChapterColor14Light}{HTML}{D8CFE6}
\definecolor{ChapterColor14Text}{HTML}{FFFFFF}
\definecolor{ChapterColor15}{HTML}{233552}      % Slate Navy
\definecolor{ChapterColor15Light}{HTML}{CFD3D9}
\definecolor{ChapterColor15Text}{HTML}{FFFFFF}
\definecolor{ChapterColor16}{HTML}{207E73}      % Sea Teal
\definecolor{ChapterColor16Light}{HTML}{CEE3E0}
\definecolor{ChapterColor16Text}{HTML}{FFFFFF}
\definecolor{ChapterColor17}{HTML}{58581D}      % Olive Brown
\definecolor{ChapterColor17Light}{HTML}{DADACD}
\definecolor{ChapterColor17Text}{HTML}{FFFFFF}

% ------------------------------------------------------------
% Code-Farben — nur für das optionale Code-Modul
% (styles/65_code_style.tex + 67_code_comments.tex) benötigt.
% Schadlos wenn das Modul in preamble.tex deaktiviert ist.
% ------------------------------------------------------------
\definecolor{ce_yellow}{rgb}{0.8,0.6,0.1}
\definecolor{ce_pink}{rgb}{0.8,0.1,0.3}
% Vordergrundfarbe des Code-Blocks. Nach ihrer Rolle benannt wie code_bg und
% code_rule — der Block ist hell, ein Token namens "white" fuer die Textfarbe
% darauf beschreibt das Gegenteil dessen, was er tut.
\definecolor{code_fg}{rgb}{0,0,0}
\definecolor{ce_green}{rgb}{0.0,0.4,0.1}
\definecolor{code_bg}{HTML}{F4F4F6}
\definecolor{code_rule}{HTML}{D8D8DE}

% ------------------------------------------------------------
% Warn- und Semantik-Farben
% ------------------------------------------------------------
\definecolor{ETHRot}{HTML}{CC0000}
\colorlet{WarnboxBack}{ETHRot!15}
% Weicher Rahmen des Warn-Tons. Eigener Wert statt des Kapitelrahmens: sonst
% ist der Warn-Ton der einzige, dessen Rahmen der Kapitelfarbe folgt — eine
% Warnbox sähe damit in jedem Kapitel anders umrandet aus.
\colorlet{WarnboxFrame}{ETHRot!25}

% Semantisches Formel-Highlighting (kapitelunabhängig, siehe \ZSFmhlA/B/C/D).
% Dedizierte Farben -- bewusst NICHT aus der Kapitelrotation (ChapterColorN), sonst
% Kollision mit der Kapitel-Identitaetsfarbe (Section 3 = Orange, Section 4 = Teal).
\definecolor{MathHLPrimary}{HTML}{00796B}    % Teal    -- Quelle / gegeben / erster Strang
\definecolor{MathHLSecondary}{HTML}{C2410C}  % Orange  -- Gegenstueck / abgeleitet / zweiter Strang
\definecolor{MathHLTertiary}{HTML}{A50D68}   % Magenta -- dritter paralleler Strang (sparsam)
\definecolor{MathHLTarget}{HTML}{1D4ED8}     % Blau    -- Ziel / Endform / Resultat

% Text-Box-Bindungen
\definecolor{TextVorBoxColor}{HTML}{0A0A1E}   % sehr dunkles Blau (einleitend)
\definecolor{TextNachBoxColor}{HTML}{0A1A0A}  % sehr dunkles Grün (ergänzend)

% Statische, kapitelunabhängige Box-Farben (Slate-Palette aehnlich Tailwind)
\definecolor{StaticBoxTitleBack}{HTML}{E2E8F0}
\definecolor{StaticBoxTitleText}{HTML}{0F172A}
\definecolor{StaticSubBoxBack}{HTML}{F8FAFC}
\definecolor{StaticSubBoxBorder}{HTML}{CBD5E1}
% Betonte Fläche des neutralen und des Warn-Tons — die Gegenstücke zu
% \FormulaboxBackColor im Kapitelton. Jede Flächen-Rolle muss in jedem Ton
% definiert sein, sonst faellt eine Kombination auf die Farbe eines fremden
% Tons zurueck (dieselbe Begruendung wie bei den drei Rahmenstaerken unten).
\colorlet{StaticEmphasisBack}{StaticBoxTitleBack!40}
\colorlet{WarnboxEmphasisBack}{ETHRot!22}
\colorlet{FormulaSeparatorColor}{black!30}
\colorlet{MutedTextColor}{black!60}
\colorlet{IdentityMarkerColor}{black!10}
\colorlet{ReleaseFooterColor}{black}
% Seitenzahl im Footer — eigener Token neben der Release-Kennung, damit beide
% Footer-Elemente unabhaengig einstellbar sind statt eines davon roh auf black.
\colorlet{ZSFPageNumberColor}{black}
% Link- und URL-Farbe. Als Rolle benannt statt als Palettenindex: "Linkfarbe"
% ist eine Entscheidung, "ChapterColor0" ist ein Slot der Kapitelrotation.
\colorlet{ZSFLinkColor}{ChapterColor0}
\colorlet{ZSFUrlColor}{ChapterColor4}
\colorlet{DangerTagColor}{red!80!black}
\newcommand{\DangerTagTextColor}{white}

% ------------------------------------------------------------
% Ink-Vertrag — wem gehoert die Farbe an dieser Stelle?
% ------------------------------------------------------------
% Farbtragende Inline-Makros (\ZSFref, \ZSFkeyword, \ZSFmhl*, die
% Groessenfarben aus 10_math) setzen ihre Farbe nicht selbst, sondern durch
% \ZSFInk. Das ist die EINE Stelle, an der entschieden wird, ob eine
% Inline-Farbe ueberhaupt angewendet wird.
%
% Der Grund ist eine stille Fehlerart: Auf einer gesaettigten Flaeche, die
% ihre eigene Kontrastfarbe setzt (Kapitel-, Abschnitts-, Unterabschnittsbalken,
% Tabellenkopf, Boxtitel, Dokumenttitel), steht eine Inline-Farbe entweder
% unlesbar auf der Flaeche oder ueberschreibt den Kontrast. Der Build meldet
% nichts, der Linter auch nicht, und die Box sieht fuer sich plausibel aus.
% Belegt an zwei realen Faellen: \ZSFref loest auf ChapterColorN auf, waehrend
% \SubsectionBarColor \chaptercolor!75 ist — derselbe Farbton bei voller
% Saettigung; und gefaerbte Mathematik im Balkentitel verliert denselben Weg.
%
% Das Kriterium ist bewusst EINES und aus dem Code ablesbar, damit es keine
% Einzelfallabwaegung braucht:
%
%   Jede Flaeche, die ihre eigene Kontrast-Textfarbe setzt, besitzt die Farbe.
%
% Wo das gilt, steht \ZSFInkOwned direkt neben dem \color/\textcolor, das den
% Kontrast setzt — nicht an einer entfernten Stelle, wo es beim naechsten
% Eingriff vergessen wird. Ist die Farbe besetzt, gibt \ZSFInk den Inhalt
% ungefaerbt aus; er erbt damit genau die Kontrastfarbe der Flaeche und bleibt
% lesbar, statt zu verschwinden.
%
% Gruppenlokal: Balken und Titel sind tcolorbox-Gruppen, der Schalter faellt
% mit ihnen zurueck. Ein manueller Reset, den man vergessen kann, existiert
% deshalb nicht.
\makeatletter
\newif\ifZSF@inkOwned
\ZSF@inkOwnedfalse
\newcommand{\ZSFInkOwned}{\ZSF@inkOwnedtrue}
\newcommand{\ZSFInkFree}{\ZSF@inkOwnedfalse}
\newcommand{\ZSFInk}[2]{{\ifZSF@inkOwned\else\color{#1}\fi#2}}
\makeatother

% ------------------------------------------------------------
% Groessenfarben — Farbe als Identitaet, nicht als Position
% ------------------------------------------------------------
% Diese Palette ist das Gegenstueck zu MathHLPrimary/Secondary/Tertiary/Target
% oben. Die vier dort sind POSITIONAL ("Quelle", "Gegenstueck", "Ziel") und
% gelten innerhalb EINER Herleitung. Diese hier sind IDENTITAET: ein Slot
% gehoert im ganzen Dokument derselben fachlichen Groesse, damit man beim
% Zusammensetzen einer Formel sieht, welche Groesse wo landet, ohne nachzulesen.
%
% Vergeben werden die Slots pro ZSF ueber \ZSFDeclareQuantity (10_math); hier
% stehen nur die Farben. Zwoelf sind bewusst die Obergrenze: Darueber hinaus
% sind Farben bei 8pt nicht mehr zuverlaessig unterscheidbar, und eine Palette,
% die man nicht auseinanderhalten kann, kostet Lesezeit statt sie zu sparen.
% Der Deklarations-Zaehler laeuft deshalb in einen Fehler und nicht in einen
% stillen Umlauf.
%
% Auswahlkriterium: als TEXT bei 8pt lesbar (also durchweg dunkel) und
% paarweise unterscheidbar — nicht "schoene Palette". Nachbarn im Farbkreis
% sind absichtlich ausgelassen.
\definecolor{ZSFQuantityColor1}{HTML}{B3261E}   % Rot
\definecolor{ZSFQuantityColor2}{HTML}{1D4ED8}   % Blau
\definecolor{ZSFQuantityColor3}{HTML}{0F766E}   % Petrol
\definecolor{ZSFQuantityColor4}{HTML}{B45309}   % Amber
\definecolor{ZSFQuantityColor5}{HTML}{6D28D9}   % Violett
\definecolor{ZSFQuantityColor6}{HTML}{15803D}   % Gruen
\definecolor{ZSFQuantityColor7}{HTML}{BE185D}   % Magenta
\definecolor{ZSFQuantityColor8}{HTML}{334155}   % Schiefer
\definecolor{ZSFQuantityColor9}{HTML}{0284C7}   % Himmelblau
\definecolor{ZSFQuantityColor10}{HTML}{7C2D12}  % Rostbraun
\definecolor{ZSFQuantityColor11}{HTML}{4D7C0F}  % Oliv
\definecolor{ZSFQuantityColor12}{HTML}{9333EA}  % Purpur
\newcommand{\ZSFQuantityPaletteSize}{12}

% \ZSFDeclareQuantity{<Name>}{<Slot 1..12>} legt \ZSFq<Name>{...} an — ein
% Makro, das seinen Inhalt in der Farbe des Slots setzt. Vergeben wird pro ZSF
% in preamble.tex, weil "welche Groessen hat dieses Fach" eine einmalige
% Entscheidung ist und keine, die pro Stelle getroffen wird:
%
%   \ZSFDeclareQuantity{Kraft}{1}      % danach: \ZSFqKraft{F}
%   \ZSFDeclareQuantity{Moment}{5}     % danach: \ZSFqMoment{M}
%
% EIN Deklarations-Makro, nicht zwei: Ob ein Name eine Groessenfamilie (Kraft,
% Moment) oder eine Rolle (zulaessig, resultierend) bezeichnet, ist eine
% fachliche Unterscheidung und keine technische — zwei Makros dafuer waeren
% derselbe Weg unter zwei Namen (rules/02_ai_mandate -> Katalog-Oekonomie).
%
% Steht hier und nicht in 10_math, obwohl das Ergebnis in Formeln landet: Die
% Entscheidung, die das Makro traegt, ist eine Farbentscheidung, und 10_math
% laedt vor diesem Modul — die Palettenpruefung unten haette dort kein
% \ZSFQuantityPaletteSize zum Pruefen. Derselbe Grund, aus dem \ZSFkeyword in
% 50_typography steht und nicht in 10_math.
%
% Zwei Fehler werden gemeldet statt still hingenommen, weil beide die
% Grundzusage brechen (eine Farbe = eine Groesse) und beide im PDF wie eine
% gueltige Einfaerbung aussehen:
%   - derselbe Slot zweimal vergeben  -> zwei Groessen waeren nicht mehr
%     unterscheidbar, obwohl die Farbe genau das verspricht;
%   - Slot ausserhalb der Palette     -> ein stiller Umlauf faerbte die erste
%     Groesse ein zweites Mal ein.
\makeatletter
\newcommand{\ZSF@quantityOwner}[1]{\csname ZSF@qslot@#1\endcsname}
\newcommand{\ZSFDeclareQuantity}[2]{%
  \ifnum#2<1\relax
    \PackageError{ZSF}{Groessen-Slot #2 liegt unter 1}%
      {Gueltig sind 1..\ZSFQuantityPaletteSize.}%
  \else
    \ifnum#2>\ZSFQuantityPaletteSize\relax
      \PackageError{ZSF}{Groessen-Slot #2 existiert nicht}%
        {Die Palette hat \ZSFQuantityPaletteSize\space Slots. Mehr Farben sind
         in der Satzgroesse der ZSF nicht mehr zuverlaessig unterscheidbar;
         eine weitere Groesse braucht deshalb keinen neuen Slot, sondern die
         Entscheidung, welche Groesse die Farbe wirklich braucht.}%
    \else
      \@ifundefined{ZSF@qslot@#2}{}{%
        \PackageError{ZSF}{Groessen-Slot #2 ist schon an
          '\ZSF@quantityOwner{#2}' vergeben}%
          {Eine Farbe gehoert genau einer Groesse — sonst sind zwei Groessen
           gleich eingefaerbt und die Farbe sagt nichts mehr aus. Einen freien
           Slot waehlen oder die vorhandene Deklaration aendern.}%
      }%
      \expandafter\gdef\csname ZSF@qslot@#2\endcsname{#1}%
      \expandafter\newcommand\csname ZSFq#1\endcsname[1]{%
        \ifZSFQuantityColors\ZSFInk{ZSFQuantityColor#2}{##1}\else##1\fi}%
    \fi
  \fi
}
\makeatother
% Globaler Hauptschalter — eine der wenigen Entscheidungen pro ZSF
% (preamble.tex). Aus, wenn die ZSF in S/W gedruckt wird oder das Fach keine
% Groessen-Identitaet braucht: die Deklarationen bleiben stehen, nur die Farbe
% entfaellt. Getrennt vom Ink-Vertrag, weil der KONTEXTBEZOGEN ist (diese
% Flaeche besitzt die Farbe) und dieser hier DOKUMENTWEIT.
\newif\ifZSFQuantityColors
\ZSFQuantityColorstrue
\newcommand{\ZSFQuantityColorsOn}{\ZSFQuantityColorstrue}
\newcommand{\ZSFQuantityColorsOff}{\ZSFQuantityColorsfalse}

% ------------------------------------------------------------
% Dokumenttitel-Masthead: bewusst unabhängig von der Kapitelrotation.
\colorlet{ZSFDocumentTitleBarColor}{ChapterColor0}
\colorlet{ZSFDocumentTitleTextColor}{ChapterColor0Text}
\colorlet{ZSFDocumentTitleBylineColor}{ChapterColor0Text!88}

% Aktive Kapitelfarben (per Kapitel-Index automatisch umgeschaltet)
% ------------------------------------------------------------
\newcommand{\chaptercolor}{ChapterColor1}
\newcommand{\chaptercolorlight}{ChapterColor1Light}
\newcommand{\chaptercolortext}{ChapterColor1Text}
\newcommand{\ChapterBarColor}{\chaptercolor}
\newcommand{\ChapterBarTextColor}{\chaptercolortext}
% Gleicher Farbton wie das Kapitel, aber eine Stufe heller: sichtbare
% Unterhierarchie bei beibehaltener weisser Schrift.
\newcommand{\SubsectionBarColor}{\chaptercolor!75}
\newcommand{\SubsectionBarTextColor}{\chaptercolortext}
\newcommand{\SubsubsectionBarColor}{\chaptercolorlight!78!white}
\newcommand{\SubsubsectionTitleColor}{\chaptercolor!72!black}
\newcommand{\FormulaboxBackColor}{\chaptercolorlight!40}
\newcommand{\TableZebraColor}{\chaptercolorlight!70}
% Skriptverweis (\ZSFScriptRef): steht IM Boxinhalt, nicht auf einem Balken.
% Deshalb der Ton für dezente Anmerkungen und nicht \ChapterBarTextColor —
% der ist die Kontrastfarbe des Balkens und damit Weiss; auf der hellen
% Boxfläche wäre der Verweis unsichtbar, ohne dass etwas meldet.
\newcommand{\ScriptRefColor}{MutedTextColor}

\newcommand{\BoxTitleColor}{\chaptercolorlight}
\newcommand{\BoxTitleTextColor}{TextVorBoxColor}
\newcommand{\BoxBackColor}{\chaptercolorlight!8}
\newcommand{\BoxFrameColor}{\chaptercolorlight!25}
% Ungetönter Boxhintergrund (tablebox, goalbox, valuegrid, leise Boxen,
% codebox). Als Token statt als literalem "white" in jeder Box-Definition:
% der Grundton aller ungetönten Boxen hängt damit an einer Zeile.
\newcommand{\BoxPlainBackColor}{white}

% Warn-Box — die einzige semantisch gefärbte Box. Ihre Farben standen bisher
% roh in der Box-Definition und waren dadurch weder auffindbar noch mit den
% übrigen Box-Tokens gemeinsam verstellbar.
\newcommand{\WarnboxBackColor}{WarnboxBack}
\newcommand{\WarnboxTitleColor}{ETHRot}
\newcommand{\WarnboxTitleTextColor}{white}

% Zielketten-Bedingung trägt bewusst eine harte Kontur statt des dezenten
% Kapitelrahmens; auch das ist eine Gestaltungsentscheidung und gehört
% deshalb hierher, nicht in styles/60_boxes.tex.
\newcommand{\GoalConditionBorderColor}{black}

% ── Ton-Tokens: wie der `tone`-Regler in den zwei Box-Familien landet ──
% Ein Ton wirkt auf Titelboxen anders als auf rahmenlose Boxen, weil die
% beiden verschiedene Träger haben: die Titelbox eine Kopfzeile und eine
% getönte Fläche, die rahmenlose Box einen Akzent und eine ruhige Fläche.
% Derselbe Regler, zwei Landepunkte — statt zweier Regler oder eines Reglers,
% der auf halber Strecke stehenbleibt (rules/02_ai_mandate).
%
% Akzent: linker Balken der leisen Box, ihre Titelzeile, das
% Aufzählungszeichen der ZSFlist.
\newcommand{\ZSFtoneAccentChapter}{\chaptercolor}
\newcommand{\ZSFtoneAccentNeutral}{StaticBoxTitleText}
\newcommand{\ZSFtoneAccentWarn}{ETHRot}

% ── Die drei Flächen-Rollen eines Tons ──
% Eine Box waehlt die ROLLE ihrer Flaeche, der Ton liefert die FARBE dazu —
% genau wie bei den drei Rahmenstaerken darunter. Deshalb muss jede Rolle in
% jedem Ton definiert sein: Eine fehlende Kombination faellt sonst auf die
% Farbe eines fremden Tons zurueck.
%
% Laute Fläche: die Füllung einer Box mit Kopfzeile.
\newcommand{\ZSFtoneLoudBackChapter}{\BoxBackColor}
\newcommand{\ZSFtoneLoudBackNeutral}{StaticSubBoxBack}
\newcommand{\ZSFtoneLoudBackWarn}{\WarnboxBackColor}

% Ruhige Fläche: die Füllung rahmenloser Boxen. Im Kapitelton bewusst
% ungetönt — die leise Box soll dezent bleiben und sich nicht der
% Titelbox angleichen. Erst ein bewusst gewählter Ton färbt sie, und genau
% dann ist die Färbung das Signal ("gehört nicht zum Kapitel" / "Warnung").
\newcommand{\ZSFtoneQuietBackChapter}{\BoxPlainBackColor}
\newcommand{\ZSFtoneQuietBackNeutral}{StaticSubBoxBack}
\newcommand{\ZSFtoneQuietBackWarn}{\WarnboxBackColor}

% Betonte Fläche: kräftiger als die laute, für Boxen, deren Inhalt selbst die
% Aussage ist (Formelbox). In jedem Ton eine Stufe über der lauten Fläche.
\newcommand{\ZSFtoneEmphasisBackChapter}{\FormulaboxBackColor}
\newcommand{\ZSFtoneEmphasisBackNeutral}{StaticEmphasisBack}
\newcommand{\ZSFtoneEmphasisBackWarn}{WarnboxEmphasisBack}

% Ungetönte Fläche (Rolle `plain`): bewusst OHNE Ton-Varianten. Sie ist die
% Abwesenheit einer Tönung und trägt in jedem Ton dieselbe Farbe — eine
% getönte Plain-Fläche wäre ein Widerspruch. Tabelle, Wertegitter, Zielkette
% und Codebox brauchen sie als neutrale Basis, auf der Zebra, Gitterlinien
% und Syntaxfarben lesbar bleiben. Der Token dafür ist \BoxPlainBackColor.

% Rahmenfarbe in drei Stärken — der Regler `frame` wählt zwischen ihnen.
% `soft` begrenzt die Box nur, `strong` grenzt sie gegen die Umgebung ab
% (Zielketten), `hard` setzt die kräftige Kontur (Formelbox). Alle drei sind
% Landepunkte desselben Tons, damit eine getönte Box in jeder Stärke stimmig
% bleibt: Jede Stärke muss in jedem Ton definiert sein, sonst fällt eine
% Kombination auf die Farbe eines fremden Tons zurück.
\newcommand{\ZSFtoneFrameSoftChapter}{\BoxFrameColor}
\newcommand{\ZSFtoneFrameSoftNeutral}{StaticSubBoxBorder}
\newcommand{\ZSFtoneFrameSoftWarn}{WarnboxFrame}
\newcommand{\ZSFtoneFrameStrongChapter}{\BoxTitleColor}
\newcommand{\ZSFtoneFrameStrongNeutral}{StaticBoxTitleText}
\newcommand{\ZSFtoneFrameStrongWarn}{ETHRot}
\newcommand{\ZSFtoneFrameHardChapter}{black}
\newcommand{\ZSFtoneFrameHardNeutral}{StaticBoxTitleText}
\newcommand{\ZSFtoneFrameHardWarn}{ETHRot}

% Für goalbox (semantische Ziel-Box)
\newcommand{\GoalConditionBackColor}{white}
\newcommand{\GoalConditionFrameColor}{\chaptercolorlight!22}
\newcommand{\GoalTargetBackColor}{\chaptercolorlight!22}
\newcommand{\GoalTargetFrameColor}{\chaptercolor}

% ------------------------------------------------------------
% Index-basierte Palette (automatisch via \StartChapter)
% ------------------------------------------------------------
\newcount\ZSFchapterPaletteIndex
\newcount\ZSFchapterPaletteLookupIndex
\newcommand{\ZSFchapterPaletteSize}{18}
\newcommand{\ZSFSetChapterPaletteByIndex}[1]{%
  \renewcommand{\chaptercolor}{ChapterColor#1}%
  \renewcommand{\chaptercolorlight}{ChapterColor#1Light}%
  \renewcommand{\chaptercolortext}{ChapterColor#1Text}%
}
\newcommand{\ZSFSetPaletteColorNameFromNumber}[2]{%
  \ZSFchapterPaletteLookupIndex=#1\relax
  \loop
  \ifnum\ZSFchapterPaletteLookupIndex>\numexpr\ZSFchapterPaletteSize-1\relax
    \advance\ZSFchapterPaletteLookupIndex by -\ZSFchapterPaletteSize
  \repeat
  \edef#2{ChapterColor\the\ZSFchapterPaletteLookupIndex}%
}
\newcommand{\SetChapterPaletteByNumber}{%
  \ZSFchapterPaletteIndex=\value{section}%
  \loop
  \ifnum\ZSFchapterPaletteIndex>\numexpr\ZSFchapterPaletteSize-1\relax
    \advance\ZSFchapterPaletteIndex by -\ZSFchapterPaletteSize
  \repeat
  \ifnum\value{section}>\numexpr\ZSFchapterPaletteSize-1\relax
    \PackageWarning{ZSF}{Chapter palette wrapped at section \thesection. Increase palette if unique colors are required.}%
  \fi
  \ZSFSetChapterPaletteByIndex{\the\ZSFchapterPaletteIndex}%
}

% ------------------------------------------------------------
% Flag-System: Spacing-Kollaps an Kapitel-/Subsection-Grenzen
% ------------------------------------------------------------
\newif\ifzsfAfterChapterBar
\zsfAfterChapterBarfalse
\newcommand{\ZSFMarkAfterChapterBar}{\global\zsfAfterChapterBartrue}
\newcommand{\ZSFConsumeAfterChapterBar}{\global\zsfAfterChapterBarfalse}

\newif\ifzsfAfterSubsectionBar
\zsfAfterSubsectionBarfalse
\newcommand{\ZSFMarkAfterSubsectionBar}{\global\zsfAfterSubsectionBartrue}
\newcommand{\ZSFConsumeAfterSubsectionBar}{\global\zsfAfterSubsectionBarfalse}

% Reduzierter "before skip" für Boxen direkt nach einer SubsectionBar.
\newcommand{\ZSFBoxBeforeSkip}{%
  \ifzsfAfterSubsectionBar\ZSFsubsectionAfterGap\else\ZSFspaceS\fi
}

% Kollabiert den before-skip der SubsectionBar auf \ZSFsubsectionAfterChapterGap,
% wenn sie direkt nach einer ChapterBar steht.
\newcommand{\ZSFCollapseAfterChapterBar}{%
  \ifzsfAfterChapterBar
    \vspace*{\dimexpr\ZSFsubsectionAfterChapterGap-\ZSFsubsectionBarBeforeSkip\relax}%
    \ZSFConsumeAfterChapterBar
  \fi
}

% ------------------------------------------------------------
% Chapter-Start-Orchestrierung
% ------------------------------------------------------------
\makeatletter
% Kurz-Wegweiser eines Front-Matter-Kapitels für das Register (siehe
% \StartFrontChapter unten). Leer heisst »dieses Kapitel ist kein Indexziel«.
% Steht hier und nicht in 66_index, weil das Index-Modul opt-in ist: Die
% Aufrufform von \StartFrontChapter darf nicht davon abhängen, ob ein Schalter
% gesetzt ist, den der Kapitelautor nicht sieht (dieselbe Begründung wie bei
% \ZSFkeyword, rules/55_index).
%
% \let auf \@empty und \def zum Setzen — NICHT \newcommand. Der Leerheitstest
% ist ein \ifx gegen \@empty, und \newcommand erzeugt ein \long-Makro, das
% \ifx nie gleich zu \@empty ist. Mit \newcommand liefe deshalb JEDER
% Registereintrag in den Front-Zweig, mit leerem Wegweiser statt einer
% Abschnittsnummer — und zwar ohne Fehler, weil beide Zweige gueltig sind.
\let\ZSF@frontIdxLabel\@empty
\newcommand{\ZSFPostChapterSpacing}{\vspace*{\ZSFpostChapterNudge}}
\newcommand{\StartChapter}[2][]{%
  \ifnum\value{zsfFirstChap}=1\setcounter{zsfFirstChap}{0}\fi
  \refstepcounter{section}
  \SetChapterPaletteByNumber
  \setcounter{subsection}{0}
  \setcounter{subsubsection}{0}
  \let\ZSF@frontIdxLabel\@empty
  \if\relax\detokenize{#1}\relax\else\label{#1}\fi
  \ChapterBar{\thesection. #2}
  \ZSFMarkAfterChapterBar
  \ZSFPostChapterSpacing
}

% Für ein Kapitel auf neuer Spalte: \ZSFNewColumn davorsetzen
% (styles/60_boxes.tex) statt einer eigenen Kapitel-Variante.

% Front-Matter-Kapitel: unnummeriert, ETH-Blau, neue Spalte ab Kapitel 2
%
%   \StartFrontChapter{Titel}              — wie bisher, kein Indexziel
%   \StartFrontChapter[Z\&E]{Zeichen …}    — mit Kurz-Wegweiser fürs Register
%
% Das optionale Argument macht ein Front-Kapitel INDEXIERBAR. Vorher war das
% ausgeschlossen, und zwar nicht durch Absicht, sondern durch Mechanik: Ein
% unnummeriertes Kapitel hat keine Abschnittsnummer, der Locator fiel auf die
% zuletzt vergebene Nummer eines fremden Kapitels zurück. Die Regel verbot
% deshalb, was das System nicht konnte — obwohl der Bedarf real ist: Ein
% Symbol- und Einheitenregister im Front-Matter ist genau die Stelle, die man
% unter Zeitdruck sucht. Statt der Nummer trägt der Eintrag jetzt diesen
% Kurznamen (Details: styles/66_index.tex).
%
% Das Label ist reiner Text und kurz — es steht in jedem Registereintrag
% dieses Kapitels und konkurriert dort mit den Abschnittsnummern.
\newcommand{\StartFrontChapter}[2][]{%
  \ifnum\value{zsfFirstChap}=1\setcounter{zsfFirstChap}{0}\else\columnbreak\fi
  \renewcommand{\chaptercolor}{ChapterColor0}%
  \renewcommand{\chaptercolorlight}{ChapterColor0Light}%
  \renewcommand{\chaptercolortext}{ChapterColor0Text}%
  \setcounter{subsection}{0}%
  \setcounter{subsubsection}{0}%
  \def\ZSF@frontIdxLabel{#1}%
  \ChapterBar{#2}%
  \ZSFMarkAfterChapterBar
  \ZSFPostChapterSpacing
}
\makeatother
